\subsection{Background Amplification in Real Data}\label{sec:res-amp}

\paragraph{Contextual gains are widely spread}
Fig.~\ref{fig:amp}(b) plots, for the most amplified field of six targets, the gain $\Delta_i(T)$ over
every background of at most two fields. The spread is large, and the maximum typically lies far
from the bulk: the South Australian wind forecast has gain $0.19$ alone, $0.20$ on average over
backgrounds, and $0.60$ over its worst background (regional available capacity and solar forecast,
which together with wind give the residual demand the gas unit PPCCGT covers). Over all 264
field--target pairs, the median ratio of the average contextual gain to the maximum is $0.32$: an
average-based score would report about a third of the worst-case risk.

\paragraph{Amplification is common}
Averaged over the 264 pairs, $\Gamma^{(1)}=0.082$, $\Gamma^{(2)}=0.110$ and $\Gamma^{(3)}=0.119$; one
background field supplies 75\% of the total amplification at $K=2$, and $K=2$ supplies 93\% of that
at $K=3$. In 91 pairs the single-field risk is below $0.3$ while the amplification exceeds $0.1$
(Fig.~\ref{fig:amp}a)---fields that look unremarkable alone and become informative in combination.
The largest amplifications are the SA wind forecast for PPCCGT ($0.19\to0.60$), the SP15
day-ahead energy price for CAISO load ($0.14\to0.51$, with the NP15 wind and SP15 solar forecasts:
price and renewables determine net load), and the SA available capacity for PPCCGT ($0.02\to0.36$).
On every pair, $\Gamma^{(1)}$ equals the strongest positive pairwise interaction
(Proposition~\ref{prop:k1}) to numerical precision.

\begin{figure}[!t]
\centering
\includegraphics[width=\textwidth]{fig_amplification.pdf}
\caption{Background amplification over ten targets. (a) Single-field risk versus amplification
$\Gamma^{(2)}$ for every field--target pair. (b) Contextual gains of the most amplified field of six
targets over all backgrounds of at most two fields, with the single-field value, the average over
backgrounds and the maximum.}
\label{fig:amp}
\end{figure}

\paragraph{Dangerous backgrounds come in families}
The top backgrounds of a field are often nearly tied and share a core partner. For the SA wind
forecast the five best backgrounds reach gains between $0.545$ and $0.603$, and all contain the SA
available capacity. We therefore report the top-$r$ backgrounds rather than a single witness; the
shared partner is the actionable information (``do not co-release with'').

\emph{The witnesses are physically meaningful.} Fig.~\ref{fig:profile} shows the field profile
of RTS line C35 and the gain matrix of NEM unit PPCCGT. The largest increases come from
combinations that reconstruct a residual quantity; for instance, the price at RTS bus 303 rises
from $0.36$ to $0.65$ for the C35 flow together with a zonal load forecast and the congestion
indicator of the neighbouring line C6 (further examples follow below). The gain matrix also shows that the same field is
harmless next to some partners and decisive next to others, which a grade that ignores the background
cannot represent.

\begin{figure}[!t]
\centering
\includegraphics[width=\textwidth]{fig4_combination_risk.pdf}
\caption{Share of candidate fields that are unsafe alone (orange), safe alone but critical with
at most two background fields (blue), or never critical (grey), at (a) $\tau=0.5$ and (b)
$\tau=0.7$; numbers give combination-critical fields over all candidates. (c) Best single field
versus best set of at most three fields.}
\label{fig:risk}
\end{figure}

\begin{figure}[!t]
\centering
\includegraphics[width=\textwidth]{fig5_profile.pdf}
\caption{(a) Field profile for the hourly flow on RTS line C35 (14 fields with the largest
$M^{(2)}$): Pearson $r^2$, single-field risk and worst-background risk, with the witness
background for the three riskiest fields. (b) Gain matrix for NEM unit PPCCGT: entry $(i,j)$ is
$V(\{i,j\})-V(\{j\})$, the diagonal is $V(\{i\})$; values below 0.10 are not printed.}
\label{fig:profile}
\end{figure}


\paragraph{Single-field grading misses most of the risk}
\begin{table}[!t]
\centering\small
\caption{Fields exposed by single-field and combination screening. $V_1$: best single field; $V_{\le3}$: best set of at most three fields; alone: fields with $V(\{i\})>\tau$; comb.: fields safe alone but $\tau$-critical with at most two background fields. Minimal unsafe sets and withholding counts are reported in Table~\ref{tab:risk-sets}.}\label{tab:risk}
\setlength{\tabcolsep}{4.5pt}
\begin{tabular}{lrrr rrrr}
\toprule
 & & & & \multicolumn{2}{c}{$\tau=0.5$} & \multicolumn{2}{c}{$\tau=0.7$}\\
\cmidrule(lr){5-6}\cmidrule(lr){7-8}
Target & $p$ & $V_1$ & $V_{\le3}$ & alone & comb. & alone & comb.\\
\midrule
RTS unit 223 & 23 & 1.00 & 1.00 & 10 & 13 & 10 & 6 \\
RTS unit 221 & 23 & 0.99 & 0.99 & 10 & 13 & 10 & 9 \\
RTS line C35 & 23 & 0.52 & 0.86 & 4 & 19 & 0 & 21 \\
NEM BW01 & 30 & 0.49 & 0.65 & 0 & 30 & 0 & 0 \\
NEM TUMUT3 & 30 & 0.38 & 0.53 & 0 & 22 & 0 & 0 \\
NEM PPCCGT & 30 & 0.55 & 0.80 & 2 & 28 & 0 & 10 \\
PJM load & 22 & 0.48 & 0.80 & 0 & 22 & 0 & 14 \\
PJM total gen. & 24 & 0.97 & 0.99 & 2 & 22 & 2 & 18 \\
PJM net interch. & 24 & 0.41 & 0.71 & 0 & 24 & 0 & 3 \\
CAISO load & 35 & 0.82 & 0.93 & 5 & 30 & 3 & 32 \\
\bottomrule
\end{tabular}
\end{table}


We first ask how much of the risk in the four data sets is invisible to single-field grading.
Table~\ref{tab:risk} and Fig.~\ref{fig:risk} report, for each of the ten targets, how many
candidate fields are unsafe on their own and how many become critical only in combination.

Of the 264 field--target
pairs, 256 are critical at $\tau=0.5$ with at most two background fields, and 223 of these are
safe on their own; at $\tau=0.7$ the figures are 138 and 113. For four targets (NEM BW01 and
TUMUT3, PJM load and PJM net interchange) no candidate exceeds $\tau=0.5$ alone, yet 22 to 30
candidates take part in a dangerous combination; at $\tau=0.7$ the same holds for RTS line C35,
where 21 of 23 fields are critical and none is unsafe alone. Withholding every field that is unsafe alone
breaks almost none of the minimal unsafe sets (Table~\ref{tab:risk-sets}): 1,766 of 1,799 at $\tau=0.5$
and 534 of 559 at $\tau=0.7$ remain.

\emph{Where single fields do leak, combinations still dominate.} Some targets are exposed by a
single close proxy: RTS unit outputs are almost determined by a nodal price ($V_1=0.99$), and
PJM total generation by the system load forecast ($0.97$). For these, $M^{(0)}$ already flags
the proxy, and the remaining risk is again combinatorial (e.g.\ 22 of 24 PJM fields are
combination-critical). Removing such near-proxies from the candidate pool
(Section~\ref{sec:data}) lowers the best single field to $0.42$ for both PJM total generation
and CAISO load but leaves every remaining field combination-critical at $\tau=0.5$ (22/22 and
26/26), so the conclusion does not depend on where the proxy boundary is drawn.
